\section{Review Design and Methods}
\subsection{Scope and analytical approach}
This is a critical integrative review with purposive selection, not a systematic review, scoping review, or meta-analysis. The modern empirical scope is January 2022 to 10 October 2026; earlier sources provide conceptual, uncertainty, and reporting context. The protocol was developed during scoping and was not preregistered. We use transparency elements consistent with sound narrative review practice \cite{sanra2019}, but do not claim PRISMA or PRISMA-ScR compliance \cite{prisma2020,prismascr2018}. The accompanying supplement specifies queries, selection rationale, extraction fields, access restrictions, and unresolved author checks.

We operationally include LLM-based healthcare systems that iterate over observations and decisions, use external tools, orchestrate components, or generate structured action proposals. Medical question answering and diagnostic dialogue are included selectively as boundary comparators: they help identify what an execution claim would require beyond a high-quality answer. Unrelated robotic control, marketing-only accounts, and unverified deployment descriptions do not provide empirical anchors for this analysis.

\subsection{Discovery and purposive selection}
Reproducible PubMed and arXiv queries combine agent terminology with clinical and healthcare terms. PubMed publication dates and arXiv submitted dates have different meanings; the scripts preserve those database-specific restrictions. Targeted retrieval from primary journal, conference, and preprint sources supplements database discovery. The original temporary exports were unavailable after an interrupted upload; the accompanying inventory is a newly retrieved snapshot, with its own timestamps. Historical counts are not presented as recovered data.

The metadata inventory is an \emph{unassessed discovery inventory}. It is not a screening flow or a denominator for publication prevalence. Identifier-level duplicate removal does not resolve article/preprint versions. Eighteen empirical reports are selected to cover diagnostic dialogue, collaborative reasoning, imaging tools, EHR tasks, management, radiotherapy planning, and administration. Eight are journal reports, four are conference reports, and six are preprints under the cited publication versions. Two are boundary comparators. These counts describe this selected map, not the literature as a whole. The absence of a particular deployment design in the map cannot establish that it does not exist elsewhere.

\subsection{Extraction and verification}
For each report we record a persistent identifier, publication status, task/sample unit, environment, tool or action type, observable effect, endpoint, reported quantitative anchor, source locator, and interpretation limit. Primary-source access is logged separately from metadata resolution. GPTPlan is extracted from its abstract only \cite{gptplan2025}; claims about its detailed implementation or patient outcomes are therefore not made. Cited versions and institutional sources take precedence over secondary news summaries.

Extraction and synthesis were AI-assisted. No independent dual screening, human consensus adjudication, agreement statistic, or formal risk-of-bias score is reported. The claim register marks independent author verification as pending. Clinical and methodological reviewers should confirm every quantitative anchor and its denominator before submission. We avoid pooling because task definitions, model configurations, tool budgets, evidence access, and scoring differ. The framework and illustrative software are original analytical outputs; they are not additional clinical efficacy evidence.
